\section{Conclusions}

In conclusion, the proposed project has helped us investigate the effectiveness of vicarious learning in the field of agentic educational learning. Even though the results did not reflect our initial hypothesis in terms of measurable knowledge gains, the Teacher and Student group consistently rated their experience higher across all learning experience dimensions, most notably in engagement and post-session confidence. This is broadly consistent with the vicarious learning hypothesis, suggesting that the presence of a synthetic peer enriches the perceived learning environment even when objective test scores remain relatively equivalent.

The main limitations are the small and informally recruited sample, the ceiling effect on simpler knowledge questions, and the audio-only delivery format. Qualitative feedback highlighted pacing and unnatural silences between turns as the most persistent issues. Future work should focus on a larger controlled study with improved TTS pacing, the incorporation of a visual teaching component, and a more varied Student agent behaviour to strengthen the pedagogical contribution of the synthetic peer.

\subsection{Future Work}

Several directions for future work were identified following the implementation and evaluation of this project.

First, the experiment should be replicated across topics of varying cognitive difficulty, as the relative simplicity of Simpson's Paradox may have contributed to the similar results observed between conditions. Second, end-of-utterance prediction should replace the current fixed response windows, which were the primary source of unnatural silences reported by participants. Third, incorporating a visual teaching component such as slide presentations would more closely replicate a conventional classroom environment. Finally, a larger and more carefully recruited sample would allow for more statistically robust conclusions.
